\documentclass[10pt,a4paper]{amsart}
\usepackage[margin=19mm]{geometry}
\usepackage{amsmath,amssymb,mathtools}
\usepackage[scr=rsfs,cal=euler]{mathalfa}
\usepackage{xfrac}
\usepackage[unicode,hidelinks,bookmarks=false]{hyperref}
\newcommand{\ie}{i.e.\ }
\newcommand{\cf}{cf.\ }
\newcommand{\resp}{respectively\ }
\newcommand{\Subsection}{\subsection}
\providecommand{\nobreakdash}{\nobreak\hbox{-}}
\setlength{\emergencystretch}{2em}
\multlinegap0pt
\usepackage{tikz}
\usetikzlibrary{cd}
\usepackage{amssymb}
\newtheorem{prop}{Proposition}
\title{On a Completion of Cohomological Functors Generalising Tate Cohomology II}
\author{Max Gheorghiu}
\date{Source: arXiv:2405.03634v3 (CC BY 4.0)}
\begin{document}
\maketitle

\noindent\emph{Source and licence.} On a Completion of Cohomological Functors Generalising Tate Cohomology II. Version arXiv:2405.03634v3 (CC BY 4.0), distributed by arXiv under the Creative Commons Attribution 4.0 International licence, http://creativecommons.org/licenses/by/4.0/, as recorded in the arXiv licence metadata for this exact version and shown on the abstract page https://arxiv.org/abs/2405.03634v3. Full licence notice: \texttt{licenses/arxiv-2405.03634-CC-BY-4.0.txt}.\par\smallskip

\noindent\textit{Editorial scope.} Counted unit: the article's prop prop:extprodsgradedcomm (source0.tex lines 1077--1097). The article's complete original proof of it is delivered with it (source0.tex lines 1099--1101, one \texttt{proof} environment of 3 lines). No mathematics is written, repaired or re-derived: the statement and the proof are the article's own lines, in the article's own notation, and no retained line is substituted or re-flowed.  Retained uncounted as the source-local context the proof needs: lines 715--720; lines 722--727; lines 745--750; lines 752--755; lines 771--774.  Nothing was omitted from any retained span, so the package carries no omission record.  No retained line contains a citation, so the wrapper carries no bibliography and no bibliography entry.  The retained lines contain the article's own diagram code, which the wrapper typesets with the standard tikz packages; no image file is delivered and the package declares no asset dependency.  Editorial formatting only, in addition: an AMS \texttt{amsart} preamble that restates the article's own declarations and macros used by the retained lines, namely the environments \texttt{prop} on separate counters, together with the standard packages those lines need.  The retained environments print in this document as Proposition 1 in order of appearance, whereas the article prints the same items with the numbering of its own full document.  The complete native source of the article as distributed in the arXiv e-print for this exact version is bundled unchanged as \texttt{sources/158737/source0.tex}.
\begin{equation}\label{diag:extandconnfirst}
\begin{tikzcd}
    \mathrm{Ext}^m(A{,} \, B) \otimes \mathrm{Ext}^n(C{,} \, E) \arrow[r, "\vee"] \arrow[d, "\delta^m \otimes \mathrm{id}"] & \mathrm{Ext}^{m+n}(A \otimes_R C{,} \, B \otimes_R E) \arrow[d, "\delta^{m+n}"] \\
    \mathrm{Ext}^{m+1}(A{,} \, B'') \otimes \mathrm{Ext}^n(C{,} \, E) \arrow[r, "\vee"] & \mathrm{Ext}^{m+n+1}(A \otimes_R C{,} \, B'' \otimes_R E)
\end{tikzcd}
\end{equation}

\begin{equation}\label{diag:extandconnsecond}
\begin{tikzcd}
    \mathrm{Ext}^m(A{,} \, B) \otimes \mathrm{Ext}^n(C{,} \, E) \arrow[r, "\vee"] \arrow[d, "\mathrm{id} \otimes \delta^n"] & \mathrm{Ext}^{m+n}(A \otimes_R C{,} \, B \otimes_R E) \arrow[d, "(-1)^m \delta^{m+n}"] \\
    \mathrm{Ext}^m(A{,} \, B) \otimes \mathrm{Ext}^{n+1}(C{,} \, E'') \arrow[r, "\vee"] & \mathrm{Ext}^{m+n+1}(A \otimes_R C{,} \, B \otimes_R E'')
\end{tikzcd}
\end{equation}

\begin{equation}\label{diag:extprodcommity}
\begin{tikzcd}
    \mathrm{Hyp}_R(A_{\bullet}{,} \, B_{\bullet})_m \otimes \mathrm{Hyp}_R(C_{\bullet}{,} \, E_{\bullet})_n \arrow[r] \arrow[d, "(-1)^{mn}\mathrm{swap}"] & \mathrm{Hyp}_R(A_{\bullet} \otimes_R C_{\bullet}{,} \, B_{\bullet} \otimes_R E_{\bullet})_{m+n} \arrow[d, "\mathrm{Hyp}_R(\mathrm{swap}_{\bullet}{,} \, \mathrm{swap}_{\bullet})_{m+n}"] \\
    \mathrm{Hyp}_R(C_{\bullet}{,} \, E_{\bullet})_n \otimes \mathrm{Hyp}_R(A_{\bullet}{,} \, B_{\bullet})_m \arrow[r] & \mathrm{Hyp}_R(C_{\bullet} \otimes_R A_{\bullet}{,} \, E_{\bullet} \otimes_R B_{\bullet})_{m+n}
\end{tikzcd}
\end{equation}

\begin{equation}\label{eq:extprodgradedcomm}
\forall x \in \mathrm{Ext}^m(A, B), y \in \mathrm{Ext}^n(C, E):
x \vee y = (-1)^{mn} y \vee x \, .
\end{equation}

\begin{prop}\label{prop:dirlimcommswithtensor}
Let $(M_i, m_i)_{i \in \mathbb{N}}$, $(N_i, n_i)_{i \in \mathbb{N}}$ be direct systems of abelian groups and denote by $\otimes$ the tensor product of abelian groups. Then
\[ \varinjlim_{i \in \mathbb{N}} (M_i \otimes N_i, m_i \otimes n_i) \cong (\varinjlim_{j \in \mathbb{N}} M_j) \otimes (\varinjlim_{k \in \mathbb{N}} N_k) \, . \]
\end{prop}

\begin{prop}\label{prop:extprodsgradedcomm}
(A form of commutativity)\newline
Assume that the tensor product $\otimes_R$ is commutative and take the homomorphism $\mathrm{swap}$ as in Diagram~\ref{diag:extprodcommity}. Then there are commutative diagrams
\begin{equation}\label{diag:swapdirlim}
\begin{tikzcd}
    {\scriptstyle \mathrm{Ext}^{m+k}(A{,} \, \widetilde{B}_k) \otimes \mathrm{Ext}^{n+k}(C{,} \, \widetilde{E}_k)} \arrow[r, "{\scriptstyle (-1)^{(m+k)(n+k)} \mathrm{swap}}"] \arrow[d, "{\scriptstyle \delta^{m+k} \otimes \mathrm{id}}"] & {\scriptstyle \mathrm{Ext}^{n+k}(C{,} \, \widetilde{E}_k) \otimes \mathrm{Ext}^{m+k}(A{,} \, \widetilde{B}_k)} \arrow[d, "{\scriptstyle (-1)^{n+k} \mathrm{id} \otimes \delta^{m+k}}"] \\
    {\scriptstyle \mathrm{Ext}^{m+k+1}(A{,} \, \widetilde{B}_{k+1}) \otimes \mathrm{Ext}^{n+k}(C{,} \, \widetilde{E}_k)} \arrow[r, shift left = 4pt, "{\scriptstyle (-1)^{(m+k+1)(n+k)} \mathrm{swap}}"] \arrow[d, "{\scriptstyle (-1)^{m+k+1} \mathrm{id} \otimes \delta^{n+k}}" near start] & {\scriptstyle \mathrm{Ext}^{n+k}(C{,} \, \widetilde{E}_k) \otimes \mathrm{Ext}^{m+k+1}(A{,} \, \widetilde{B}_{k+1})} \arrow[d, "{\scriptstyle \delta^{n+k} \otimes \mathrm{id}}"] \\
    {\scriptstyle \mathrm{Ext}^{m+k+1}(A{,} \, \widetilde{B}_{k+1}) \otimes \mathrm{Ext}^{n+k+1}(C{,} \, \widetilde{E}_{k+1})} \arrow[r, shift left = 4pt, "{\scriptstyle (-1)^{(m+k+1)(n+k+1)} \mathrm{swap}}"] & {\scriptstyle \mathrm{Ext}^{n+k+1}(C{,} \, \widetilde{E}_{k+1}) \otimes \mathrm{Ext}_R^{m+k+1}(A{,} \, \widetilde{B}_{k+1})}
\end{tikzcd}
\end{equation}
The homomorphism in their direct limit
\[ \widehat{\mathrm{swap}}: \widehat{\mathrm{Ext}}^m(A, B) \otimes \widehat{\mathrm{Ext}}^n(C, E) \rightarrow \widehat{\mathrm{Ext}}^n(C, E) \otimes \widehat{\mathrm{Ext}}^m(A, B) \]
renders external products commutative in the sense that the diagram
\begin{equation}\label{diag:fakecommity}
\begin{tikzcd}
    \widehat{\mathrm{Ext}}^m(A{,} \, B) \otimes \widehat{\mathrm{Ext}}^n(C{,} \, E) \arrow[r, "\vee"] \arrow[d, "\widehat{\mathrm{swap}}"] & \widehat{\mathrm{Ext}}^{m+n}(A \otimes_R C{,} \, B \otimes_R E) \arrow[d, "\widehat{\mathrm{Ext}}^{m+n}(\mathrm{swap}{,} \, \mathrm{swap})"] \\
    \widehat{\mathrm{Ext}}^n(C{,} \, E) \otimes \widehat{\mathrm{Ext}}^m(A{,} \, B) \arrow[r, "\vee"] & \widehat{\mathrm{Ext}}^{m+n}(B \otimes_R E{,} \, A \otimes_R C)
\end{tikzcd}
\end{equation}
commutes. Cup products of complete cohomology satisfy the same form of commutativity.
\end{prop}

\begin{proof}
Diagram~\ref{diag:swapdirlim} commutes by definition. If we apply Proposition~\ref{prop:dirlimcommswithtensor}, the homomorphisms on the left hand side of Diagram~\ref{diag:swapdirlim} do not immediately yield $\widehat{\mathrm{Ext}}^m(A, B) \otimes \widehat{\mathrm{Ext}}^n(C, E)$ in the direct limit as in resolution construction. However, if we concatenate four copies of Diagram~\ref{diag:swapdirlim}, the sign issues vanish yielding the direct limit of the desired form. For the same reason the homomorphisms on the right hand side result in the correct direct limit from which we obtain the homomorphism $\widehat{\mathrm{swap}}$. Next, we form with Diagram~\ref{diag:extprodcommity} the following direct system of commuting squares. We connect its left hand homomorphisms via Diagram~\ref{diag:swapdirlim} to each other. We connect its top and its bottom homomorphisms via Diagram~\ref{diag:extandconnfirst} and Diagram~\ref{diag:extandconnsecond} where the factor of $(-1)^m$ is assigned to the term $\mathrm{id} \otimes \delta^n$ instead of the term $\delta^{m+n}$ in the latter diagram. According to Equation~\ref{eq:extprodgradedcomm}, the right hand morphisms are of the form $\mathrm{Ext}^{m+n+2k}(\mathrm{swap}, \mathrm{swap})$ and connected by connecting homomorphism. Then Diagram~\ref{diag:fakecommity} results as the direct limit of these squares.
\end{proof}

\end{document}
